% ═══════════════════════════════════════════════════════════════
% SPANISCH ARBEITSBLATT AB-03b: hay + muebles
% ═══════════════════════════════════════════════════════════════
% Sequenz 3: Mi casa | Block b | Klassenstufe 7
% Gesamtpunktzahl: 20 Punkte
% ═══════════════════════════════════════════════════════════════

\documentclass[11pt, a4paper]{article}

% ═══════════════════════════════════════════════════════════════
% SW-MODUS TOGGLE
% ═══════════════════════════════════════════════════════════════
\newif\ifswmodus
\swmodusfalse

% ═══════════════════════════════════════════════════════════════
% PAKETE
% ═══════════════════════════════════════════════════════════════

\usepackage[utf8]{inputenc}
\usepackage[T1]{fontenc}
\usepackage[ngerman]{babel}
\usepackage{helvet}
\renewcommand{\familydefault}{\sfdefault}

\usepackage[margin=2cm]{geometry}
\usepackage{enumitem}
\usepackage{tabularx}
\usepackage{booktabs}
\usepackage{multirow}
\usepackage{graphicx}

\usepackage{xcolor}
\usepackage{framed}
\usepackage{tcolorbox}
\tcbuselibrary{skins}

\usepackage{fancyhdr}
\usepackage{lastpage}
\usepackage{needspace}
\usepackage{tikz}

% ═══════════════════════════════════════════════════════════════
% FARBEN (sciencesim-konform)
% ═══════════════════════════════════════════════════════════════

\definecolor{spanisch-color}{HTML}{c41e3a}
\definecolor{grau-color}{HTML}{888888}
\definecolor{hellgrau-color}{HTML}{f5f5f5}
\definecolor{orange-color}{HTML}{b35c00}

\definecolor{sw-dark}{HTML}{000000}
\definecolor{sw-gray}{HTML}{666666}
\definecolor{sw-white}{HTML}{ffffff}

\ifswmodus
    \colorlet{spanisch}{sw-dark}
    \colorlet{grau}{sw-gray}
    \colorlet{hellgrau}{sw-white}
    \colorlet{orange}{sw-dark}
\else
    \colorlet{spanisch}{spanisch-color}
    \colorlet{grau}{grau-color}
    \colorlet{hellgrau}{hellgrau-color}
    \colorlet{orange}{orange-color}
\fi

\newcommand{\fachfarbe}{spanisch}

% ═══════════════════════════════════════════════════════════════
% VARIABLEN
% ═══════════════════════════════════════════════════════════════

\newcommand{\thema}{hay + muebles (Es gibt + Möbel)}
\newcommand{\sequenz}{03}
\newcommand{\block}{b}
\newcommand{\fach}{Spanisch}
\newcommand{\klasse}{7}
\newcommand{\maxpunkte}{20}
\newcommand{\datum}{\today}

% ═══════════════════════════════════════════════════════════════
% KOPF- UND FUSSZEILE
% ═══════════════════════════════════════════════════════════════

\pagestyle{fancy}
\fancyhf{}
\renewcommand{\headrulewidth}{0.4pt}
\renewcommand{\footrulewidth}{0.4pt}

\lhead{\fach{} -- Klasse \klasse}
\chead{AB-\sequenz\block}
\rhead{\datum}

\lfoot{}
\cfoot{}
\rfoot{Seite \thepage\ von \pageref{LastPage}}

% ═══════════════════════════════════════════════════════════════
% BEFEHLE
% ═══════════════════════════════════════════════════════════════

\newcommand{\punktebox}[1]{\hfill\fbox{\small \_/#1}}

\newcommand{\luecke}[1]{\rule{#1}{0.4pt}}

\newtcolorbox{merkkasten}[1][]{
    colback=hellgrau,
    colframe=\fachfarbe,
    colbacktitle=white,
    coltitle=\fachfarbe,
    boxrule=1pt,
    arc=0pt,
    left=8pt, right=8pt, top=8pt, bottom=8pt,
    fonttitle=\bfseries,
    attach boxed title to top left={yshift=-2mm, xshift=5mm},
    boxed title style={boxrule=0pt, colback=white},
    title=#1
}

\newtcolorbox{vokabelkasten}{
    colback=white,
    colframe=\fachfarbe,
    boxrule=0.5pt,
    arc=0pt,
    left=5pt, right=5pt, top=5pt, bottom=5pt
}

\newtcolorbox{hinweis}{
    colback=white,
    colframe=orange,
    colbacktitle=white,
    coltitle=orange,
    boxrule=1pt,
    arc=0pt,
    left=5pt, right=5pt, top=5pt, bottom=5pt,
    fonttitle=\bfseries,
    attach boxed title to top left={yshift=-2mm, xshift=5mm},
    boxed title style={boxrule=0pt, colback=white},
    title=Hinweis
}

\newenvironment{aufgabe}[1]{%
    \needspace{5\baselineskip}%
    \nopagebreak[4]%
    \vspace{10pt}
    \noindent\textbf{\textcolor{\fachfarbe}{Aufgabe #1}}
    \nopagebreak[4]%
    \vspace{5pt}
    \nopagebreak[4]%
    \begin{list}{}{\leftmargin=0pt}
    \item[]
}{%
    \end{list}
}

\newlist{teilaufgaben}{enumerate}{2}
\setlist[teilaufgaben,1]{label=\alph*), leftmargin=2em}
\setlist[teilaufgaben,2]{label=\roman*), leftmargin=2em}

\newcommand{\es}[1]{\textcolor{\fachfarbe}{\textbf{#1}}}

% ═══════════════════════════════════════════════════════════════
% DOKUMENT
% ═══════════════════════════════════════════════════════════════

\begin{document}

% ═══════════════════════════════════════════════════════════════
% KOPFBEREICH
% ═══════════════════════════════════════════════════════════════

\begin{center}
    {\Large\bfseries\textcolor{\fachfarbe}{Arbeitsblatt: \thema}}
\end{center}

\vspace{5pt}

\noindent
\begin{tabularx}{\textwidth}{|X|X|X|}
\hline
\textbf{Name:} \luecke{4cm} &
\textbf{Klasse:} \klasse &
\textbf{Datum:} \luecke{2cm} \\
\hline
\end{tabularx}

\vspace{15pt}

% ═══════════════════════════════════════════════════════════════
% MERKKASTEN 1: hay
% ═══════════════════════════════════════════════════════════════

\begin{merkkasten}[Merkkasten: hay (es gibt)]
\textbf{Die Struktur \es{hay}:}

\vspace{0.5em}

\es{Hay} bedeutet ``\luecke{3cm}'' und wird für Singular \textbf{und} Plural verwendet.

\vspace{0.8em}

\textbf{Struktur:}

\begin{center}
\fbox{\es{hay} + \es{un/una} + Substantiv}
\end{center}

\vspace{0.5em}

\textbf{Beispiele:}

\begin{tabular}{ll}
$\bullet$ En mi dormitorio \es{hay una} cama. & = \luecke{5cm} \\[0.6em]
$\bullet$ En el salón \es{hay un} sofá. & = \luecke{5cm} \\
\end{tabular}

\vspace{0.8em}

\textbf{ACHTUNG:} Nach \es{hay} steht der \textbf{un}bestimmte Artikel (un/una), \textbf{nicht} el/la!

\end{merkkasten}

\vspace{10pt}

% ═══════════════════════════════════════════════════════════════
% MERKKASTEN 2: Möbelvokabeln
% ═══════════════════════════════════════════════════════════════

\begin{merkkasten}[Merkkasten: Los muebles (Die Möbel)]

\vspace{0.5em}

\begin{tabular}{ll|ll}
\es{la mesa} & = \luecke{3cm} & \es{el armario} & = \luecke{3cm} \\[0.8em]
\es{la silla} & = \luecke{3cm} & \es{la lámpara} & = \luecke{3cm} \\[0.8em]
\es{el sofá} & = \luecke{3cm} & \es{la estantería} & = \luecke{3cm} \\[0.8em]
\es{la cama} & = \luecke{3cm} & \es{el escritorio} & = \luecke{3cm} \\
\end{tabular}

\vspace{0.5em}

\textit{Tipp: Lerne die Wörter immer mit Artikel! (el = maskulin, la = feminin)}

\end{merkkasten}

\vspace{15pt}

% ═══════════════════════════════════════════════════════════════
% AUFGABE 1: Möbel zuordnen
% ═══════════════════════════════════════════════════════════════

\begin{aufgabe}{1: Möbel zuordnen}
Verbinde das spanische Wort mit der deutschen Bedeutung. \punktebox{4}
\end{aufgabe}

\begin{center}
\begin{tabular}{rcl}
\es{la cama} & \luecke{3cm} & der Schrank \\[0.8em]
\es{el escritorio} & \luecke{3cm} & das Bett \\[0.8em]
\es{la estantería} & \luecke{3cm} & der Schreibtisch \\[0.8em]
\es{el armario} & \luecke{3cm} & das Regal \\
\end{tabular}
\end{center}

\vspace{15pt}

% ═══════════════════════════════════════════════════════════════
% AUFGABE 2: hay + un/una einsetzen
% ═══════════════════════════════════════════════════════════════

\begin{aufgabe}{2: hay + un/una einsetzen}
Ergänze die Sätze mit \es{hay un} oder \es{hay una}. \punktebox{4}
\end{aufgabe}

\begin{vokabelkasten}
\textit{Zur Erinnerung:} un = maskulin (el) | una = feminin (la)
\end{vokabelkasten}

\vspace{0.5em}

\noindent
a) En mi dormitorio \luecke{3cm} cama grande.

\vspace{0.6em}
\noindent
b) En el salón \luecke{3cm} sofá nuevo.

\vspace{0.6em}
\noindent
c) En la cocina \luecke{3cm} mesa.

\vspace{0.6em}
\noindent
d) En mi habitación \luecke{3cm} escritorio.

\vspace{15pt}

% ═══════════════════════════════════════════════════════════════
% AUFGABE 3: Zimmer beschreiben
% ═══════════════════════════════════════════════════════════════

\begin{aufgabe}{3: Zimmer beschreiben}
Was gibt es in diesen Zimmern? Schreibe Sätze mit ``En... hay...''. \punktebox{6}
\end{aufgabe}

\begin{hinweis}
Beginne jeden Satz mit dem Zimmer: \es{En el dormitorio}, \es{En el salón}, \es{En la cocina}...
\end{hinweis}

\vspace{0.8em}

\noindent
\textbf{Beispiel:} \textit{Schlafzimmer: Bett, Schrank}

$\rightarrow$ \es{En el dormitorio hay una cama y un armario.}

\vspace{1em}

\noindent
a) \textbf{Wohnzimmer:} Sofa, Lampe

\vspace{0.5em}
\noindent
\rule{\textwidth}{0.4pt}

\vspace{1em}

\noindent
b) \textbf{Küche:} Tisch, 4 Stühle

\vspace{0.5em}
\noindent
\rule{\textwidth}{0.4pt}

\vspace{1em}

\noindent
c) \textbf{Arbeitszimmer:} Schreibtisch, Regal, Lampe

\vspace{0.5em}
\noindent
\rule{\textwidth}{0.4pt}

\vspace{15pt}

% ═══════════════════════════════════════════════════════════════
% AUFGABE 4: Eigenes Zimmer beschreiben
% ═══════════════════════════════════════════════════════════════

\begin{aufgabe}{4: Dein Zimmer beschreiben}
Beschreibe dein eigenes Zimmer auf Spanisch. Schreibe mindestens 4 Sätze mit \es{hay}. \punktebox{6}
\end{aufgabe}

\begin{hinweis}
Du kannst auch schreiben, was es \textbf{nicht} gibt: \es{No hay...} (Es gibt kein/e...)
\end{hinweis}

\vspace{0.8em}

\noindent
\textbf{Mi habitación / Mi dormitorio:}

\vspace{0.8em}
\noindent
\rule{\textwidth}{0.4pt}

\vspace{0.8em}
\noindent
\rule{\textwidth}{0.4pt}

\vspace{0.8em}
\noindent
\rule{\textwidth}{0.4pt}

\vspace{0.8em}
\noindent
\rule{\textwidth}{0.4pt}

\vspace{0.8em}
\noindent
\rule{\textwidth}{0.4pt}

\vspace{0.8em}
\noindent
\rule{\textwidth}{0.4pt}

% ═══════════════════════════════════════════════════════════════
% PUNKTEÜBERSICHT
% ═══════════════════════════════════════════════════════════════

\vspace{20pt}
\hrule
\vspace{10pt}

\begin{flushright}
\textbf{Punkte:} \luecke{1cm} / \maxpunkte
\end{flushright}

\end{document}
